\chapter{Heartburn}

\section*{Definition}
Heartburn is a burning retrosternal discomfort or pain that rises from the epigastrum toward the throat, typically occurring after meals or when lying down. It results from gastroesophageal reflux of acidic gastric contents into the esophagus and represents the most common symptom of gastroesophageal reflux disease (GERD).

\section*{What Happens in Your Body}
Heartburn occurs when the lower esophageal sphincter (LES) is incompetent or relaxes inappropriately, allowing gastric acid, bile, and pepsin to reflux into the distal esophagus. The esophageal mucosa lacks the protective barrier of gastric epithelium, making it vulnerable to injury. Transient LES relaxations (tLESRs) are the primary mechanism in most people. Risk factors include obesity, hiatal hernia, delayed gastric emptying, and increased intra-abdominal pressure. Chronic exposure leads to esophagitis, strictures, Barrett esophagus, and esophageal adenocarcinoma.

\section*{Causes}
\begin{itemize}
  \item Gastroesophageal reflux disease (GERD) — most common cause
  \item Hiatal hernia (sliding or paraesophageal)
  \item Functional heartburn (no visible mucosal injury on endoscopy)
  \item Eosinophilic esophagitis (allergic inflammatory condition)
  \item Infectious esophagitis (Candida, HSV, CMV — typically in immunocompromised)
  \item Pill esophagitis (NSAIDs, bisphosphonates, potassium, tetracyclines, iron)
  \item Gastroparesis (delayed gastric emptying, common in diabetes)
  \item Pregnancy (hormonal relaxation of LES and increased abdominal pressure)
  \item Obesity (increased intra-abdominal pressure promoting reflux)
  \item Lifestyle triggers (large meals, fatty/spicy foods, caffeine, alcohol, tobacco, chocolate, mint)
\end{itemize}

\section*{When to See a Doctor}
See your doctor if:
\begin{itemize}
  \item Heartburn occurring more than twice per week
  \item Symptoms persisting despite over-the-counter antacids or lifestyle modifications
  \item Difficulty or pain with swallowing (dysphagia, odynophagia)
  \item Unexplained weight loss, persistent nausea/vomiting, or signs of GI bleeding
  \item Chest pain with exertion or pressure (must exclude cardiac causes first)
\end{itemize}

\section*{Related Symptoms}
\begin{itemize}
  \item Regurgitation of sour or bitter fluid into the mouth
  \item Dysphagia (difficulty swallowing) or odynophagia (painful swallowing)
  \item Chronic cough or hoarseness (laryngopharyngeal reflux)
  \item Globus sensation (lump in throat)
  \item Belching or bloating
  \item Asthma exacerbations (reflux-triggered bronchospasm)
  \item Nausea and epigastric fullness
  \item Dental erosions from acid exposure
\end{itemize}

\section*{Self-Care / Home Management}
\begin{itemize}
  \item Eat smaller, more frequent meals and avoid lying down for at least 2-3 hours after eating.
  \item Elevate the head of the bed by 6-8 inches using blocks or a wedge pillow.
  \item Avoid trigger foods: fatty/fried foods, chocolate, mint, caffeine, alcohol, spicy foods, citrus, and tomatoes.
  \item Maintain a healthy weight, stop smoking, and avoid tight clothing around the abdomen.
  \item Over-the-counter antacids (calcium carbonate, magnesium/aluminum hydroxide) or H2 blockers provide short-term relief.
\end{itemize}

\section*{How Your Doctor Will Evaluate You}
\begin{itemize}
  \item History covers onset, duration, frequency, severity, aggravating/relieving factors, and red flag symptoms (dysphagia, weight loss, GI bleeding).
  \item Upper endoscopy (esophagogastroduodenoscopy, EGD) is the gold standard for evaluating GERD complications, esophagitis, Barrett esophagus, and excluding malignancy.
  \item Ambulatory pH monitoring measures esophageal acid exposure time and correlates symptoms with reflux episodes.
  \item Esophageal manometry evaluates LES pressure and esophageal peristalsis, particularly before antireflux surgery.
\end{itemize}

\section*{Treatment Options}
\begin{itemize}
  \item First-line therapy includes lifestyle modifications and acid suppression with proton pump inhibitors (PPIs) or H2 receptor antagonists.
  \item Prokinetic agents may be added for people with delayed gastric emptying or weak peristalsis.
  \item Antireflux surgery (Nissen fundoplication) is reserved for well-selected people with confirmed GERD who are intolerant of or refractory to medical therapy.
  \item Endoscopic therapies (Stretta procedure, transoral incisionless fundoplication) are less invasive alternatives for appropriate candidates.
\end{itemize}

\clearpage
